% !TeX spellcheck = pt_BR

\chapter{Sobre a fronteira entre as equações experimentais e demonstráveis} \label{apendice-eqs-exp-vs-dem}



À primeira vista, a diferença entre uma equação experimental e uma
demonstrável parece totalmente óbvia, e até mesmo definitiva. Uma equação experimental, como mostramos, nasce
de medições: o cientista observa o mundo, identifica padrões e os sintetiza
em linguagem matemática. Uma equação demonstrável, por outro lado, nasce de deduções: por meio da matemática, partindo de princípios já estabelecidos e, por desenvolvimento
lógico, chega-se a uma nova conclusão. A fronteira entre as duas, contudo,
é menos nítida do que parece --- e entender por que é assim mostra algo
profundo não apenas sobre a Física, mas sobre a própria natureza do
conhecimento científico. O autor já deixa de sobreaviso que este se trata de um apêndice um tanto filosófico e busca justificar para o leitor mais curioso algumas nuances sobre os tipos de equações físicas e a classificação usada pelo autor nessa obra. Avisos feitos, sigamos para a nossa reflexão.


\section*{A teoria que explica o experimento}

Considere o seguinte cenário: um cientista realiza experimentos
meticulosos, coleta dados durante anos e, ao final, sintetiza tudo em uma
equação que descreve fielmente o que observou. Trata-se,
inequivocamente, de uma equação experimental. Décadas depois, outro
cientista constrói uma teoria mais ampla --- uma estrutura de princípios
e definições --- e demonstra que aquela equação experimental é uma
consequência lógica necessária dessa teoria. O que era experimental
torna-se, agora, demonstrável.

O que aconteceu com o estatuto epistemológico\footnote{A epistemologia é o ramo da filosofia, também conhecido como Teoria do Conhecimento, que estuda a natureza, origem, métodos e limites do conhecimento humano. Ela investiga como adquirimos saberes, o que torna uma crença verdadeira e justificável, e diferencia o conhecimento científico do senso comum. O aviso foi dado de que essa era uma discussão um tanto quanto filosófica.} da equação? Ela mudou de
tipo? Ou ela sempre foi as duas coisas ao mesmo tempo, dependendo apenas
do ponto de vista de quem a olha? A resposta honesta é: depende do
referencial histórico e teórico que se adota. Assim como o movimento, o tipo de categoria que uma equação assume na nossa classificação é, em alguma medida, relativo. Uma equação não carrega
seu tipo inscrito na tinta -- ele é, em certa medida, relativo ao
estado do conhecimento disponível.



Não é preciso ir longe para encontrar o melhor exemplo dessa
ambiguidade: ele está no capítulo de Gravitação deste livro. Johannes Kepler, no início do século~XVII, passou anos debruçado sobre
as observações astronômicas de Tycho Brahe e delas extraiu três leis que
descrevem com extraordinária precisão o movimento dos planetas. A
Terceira Lei,

$$ \frac{T^2}{R^3} = \text{constante}, $$

é um resultado empírico puro: Kepler não sabia por que os planetas
obedeciam a essa relação. Ele tinha a descrição; faltava a explicação. 
É por isso que a classificamos, neste livro, como \textbf{Experimental} no seu estudo --- trata-se da Equação \eqref{eq_gravitacao_terceiraleikepler}.

Isaac Newton chegou meio século depois com a explicação. Ao postular a
Lei da Gravitação Universal e combiná-la com a Segunda Lei de Newton e a
condição de movimento circular uniforme, obteve-se a Terceira Lei de
Kepler como consequência algébrica direta, como fizemos no estudo da Equação \eqref{eq_gravitacao_leigravitacaouniversal}.

O que Kepler havia extraído de tabelas de observação, Newton demonstrou --- partindo de outra lei --- em algumas linhas de álgebra. A Terceira Lei de Kepler é, portanto,
simultaneamente experimental --- no sentido histórico, no sentido de
como ela foi obtida pela primeira vez --- e demonstrável --- no sentido
lógico, dentro da teoria newtoniana.

Há ainda um detalhe a mais nessa história, o qual foi discutido por nós no capítulo de gravitação, mas merece que venha à tona neste momento. A própria teoria de Newton
foi \textit{construída} para ser compatível com as leis de Kepler. A
dependência $1/r^2$ da força gravitacional não foi obtida ao acaso: ela foi deduzida a partir da Terceira Lei de Kepler,
usando o argumento centrípeto, como mostramos. Newton não descobriu a gravitação e
depois verificou que ela explicava as leis de Kepler. Ele construiu a gravitação
\textit{de modo que} ela explicasse Kepler.\footnote{De bandeja, a teoria acabou explicando vários outros aspectos do movimento planetário que transcendiam o que as leis de Kepler eram capazes de prever.} Isso cria uma circularidade
lógica: a lei experimental de Kepler é usada para
construir a teoria de Newton, que por sua vez, é capaz de demonstrar a lei de Kepler.
As duas equações carregam, em certo sentido, o mesmo conteúdo físico
--- apenas empacotado de formas diferentes.


\section*{Organizando o caos por meio da razão}

Há uma forma de enxergar essa relação entre experimento e teoria que vai
além da Física e toca a filosofia do conhecimento. Em sua
\textit{Crítica da Razão Pura}, Kant argumenta que o mundo em si --- o
mundo tal como ele é, independentemente de nós --- é, em certa medida,
caótico e informe. O que fazemos, como sujeitos racionais, é
\textit{impor ordem} a esse caos: nossa razão projeta sobre a
experiência bruta categorias, estruturas e princípios que tornam o mundo
inteligível.

Guardadas as devidas proporções e sem forçar uma leitura kantiana
estrita da Física, essa ideia captura algo verdadeiro sobre o que Newton
fez com as observações de Kepler. Os dados brutos de Tycho Brahe e as leis
indutivas de Kepler eram, em certo sentido, o caos: padrões numéricos
precisos, mas sem explicação, sem estrutura unificadora, sem razão de
ser. Newton impôs ordem sobre esse caos ao construir uma teoria --- um
sistema lógico de princípios --- que não apenas descrevia os padrões, mas os
tornava \textit{necessários}. Dentro da teoria newtoniana, os planetas
\textit{têm que} obedecer à Terceira Lei de Kepler: não é uma
coincidência, é uma consequência.

Toda grande teoria física faz isso: ela substitui um conjunto de fatos
experimentais dispersos por uma estrutura lógica unificada da qual esses
fatos emergem como conclusões inevitáveis --- meras consequências. O experimento revela
\textit{o que} acontece; a teoria explica \textit{por que} não poderia
ser diferente. E quando a teoria chega, as equações experimentais que a
motivaram ganham um novo estatuto: elas se tornam demonstráveis, sem
deixar de terem sido experimentais.


\section*{A categoria oculta no livro: os princípios fundamentais}

Por fim, cabe um comentário sobre um quarto tipo de equação que a classificação do GFE, deliberadamente, não nomeia --- mas que o leitor atento certamente já percebeu.

Considere as Leis de Newton. A Primeira Lei não é uma definição: ela não
cria nenhum conceito novo. Tampouco é experimental no sentido de ter
sido extraída de dados: nenhum experimento pode confirmar diretamente
que um corpo \textit{sem força alguma} se move em linha reta com
velocidade constante, pois é impossível eliminar completamente todas as
forças em um experimento real, como vimos no experimento dos planos inclinados de Galileu: a conclusão é um \emph{passo diante do abismo,} uma \emph{abstração.} Assim, tal conteúdo certamente não é demonstrável: não existe nenhum princípio mais fundamental dentro da mecânica newtoniana a partir
do qual ela se derive. Ela é, simplesmente, um \textit{postulado} --- um
axioma do sistema teórico, aceito como ponto de partida porque o
edifício que se constrói sobre ele funciona extraordinariamente bem.

O mesmo vale para a Terceira Lei de Newton. A ação e a reação não se
deduzem das outras duas leis: trata-se de um princípio independente, adicionado
ao sistema porque a natureza, ao que tudo indica, assim se comporta. E
vale, em certa medida, para a própria Segunda Lei: classificamo-la neste
livro como \textbf{Definição} --- e essa é uma posição defensável, pois
ela define operacionalmente o conceito de força. Mas há quem argumente,
com razão, que $F = ma$ é menos uma definição livre e mais um postulado
estrutural: a afirmação de que existe uma grandeza chamada força,
mensurável independentemente de seus efeitos, que se relaciona com massa
e aceleração precisamente dessa forma. A distinção é sutil, mas existe.

Esse quarto tipo de equação poderia ser chamado de \textbf{Princípio} --- ou Axioma, ou Postulado. Ele se distingue das Definições porque não é
arbitrário: não poderíamos escolher que $F = m^2a$ e construir uma
mecânica igualmente coerente\footnote{É fácil de perceber isso: se a equação fundamental da dinâmica (Segunda Lei de Newton) fosse $F=m^2a,$ um corpo caindo sob a ação da gravidade (força peso) teria como força resultante $F_R = P =mg$, ou seja,  $  m^{2} a =  mg$, e a aceleração de corpos em queda seria dada por $a=g/m,$ de modo que ela dependeria da massa do corpo em queda, contradizendo os fatos experimentais.}. Distingue-se das Experimentais porque não
emerge diretamente de um experimento identificável. E distingue-se das
Demonstráveis porque não se deriva de nada mais fundamental --- ele
\textit{é} o mais fundamental.


Por que o GFE não criou essa quarta categoria? Por uma razão pedagógica. A distinção entre Definição, Experimental e Demonstrável já é
suficientemente não trivial para um primeiro contato --- ela exige que
o leitor pense sobre a origem de cada equação --- dentro do contexto histórico em que ela surge --- algo que
a maioria dos livros didáticos simplesmente ignora. Acrescentar uma
quarta categoria, cujas fronteiras com as demais são ainda mais sutis e
filosóficas, arriscaria sobrecarregar o leitor no momento em que ele
ainda está aprendendo a usar o sistema. A opção foi simplificar: as
Leis de Newton foram absorvidas nas categorias que o autor julgou serem as mais próximas --- a
Segunda como Definição, as demais como Experimentais --- e a riqueza
epistemológica completa foi relegada a este apêndice, que o leitor mais
curioso, espera-se, encontrará no momento certo.



\section*{A escolha metodológica deste livro}

Diante de toda essa ambiguidade, cabe destacar uma escolha que foi feita na composição dessa obra. A opção adotada foi a
seguinte: classificar cada equação de acordo com o seu estatuto
\textit{no momento em que ela aparece no texto}, segundo o
desenvolvimento cronológico e lógico escolhido.

A Terceira Lei de Kepler aparece \textit{antes} da Lei da Gravitação
Universal. No momento em que é apresentada, não há nenhuma teoria
disponível para demonstrá-la --- ela é, naquele ponto da narrativa,
genuinamente experimental. Por isso, é classificada como Experimental.
Que ela possa ser demonstrada a partir da gravitação newtoniana é uma
informação que o leitor só terá algumas páginas depois, e está
explicitada na seção de Observações da equação gravitacional.

O autor acredita que essa escolha tem uma importante virtude pedagógica: ela preserva a ordem
histórica e epistêmica do descobrimento. Apresentar a Terceira Lei de
Kepler como demonstrável, logo de saída, seria um anacronismo --- seria
contar o final da história antes do começo. O leitor que acompanha o
texto na ordem cronológica vivencia, em miniatura, o mesmo arco
intelectual que a humanidade percorreu: primeiro a observação bruta,
depois a teoria que a unifica e a explica. Primeiro o \emph{caos} de Kepler;
depois a \emph{ordem} de Newton.

O leitor maduro deve ter em mente que os três tipos de classificações das equações do GFE ---
Definição, Experimental e Demonstrável --- são ferramentas pedagógicas,
não categorias ontológicas absolutas. Elas organizam o conhecimento de
uma forma que o autor julgou mais clara e honesta; não pretendem ser a
última palavra sobre a natureza das equações físicas. O que elas
pretendem, isso sim, é que o leitor nunca perca de vista a pergunta
mais importante de todas: \textit{de onde vem o direito de escrever
essa equação?} Ou seja, de onde surgem as \emph{fórmulas} e por que elas fazem sentido? Às vezes, a resposta é: \textit{medimos}. Às vezes, é:
\textit{deduzimos}. Às vezes, é: \textit{postulamos, e o mundo
concordou}. Às vezes, é: \textit{medimos, e então construímos uma teoria
que tornou a medição necessária}. Essa última resposta é, talvez, a
mais bela das quatro --- e é exatamente ela que descreve o que Kepler
e Newton fizeram juntos, ainda que separados um do outro por meio século.